\documentclass[12pt]{article}
\usepackage{graphicx} % Required for inserting images
\usepackage[margin=0.5in]{geometry}
\usepackage{amsmath, amssymb}
\usepackage{mathrsfs}


\usepackage{soul}


% Dr. David Brown Commands
\newcommand{\ds}{\displaystyle}
\newcommand{\R}{\mathbb{R}}
\newcommand{\N}{\mathbb{N}}
\newcommand{\Q}{\mathbb{Q}}
\newcommand{\C}{\mathbb{C}}


% Commands to get script r
\usepackage{calligra}
\DeclareMathAlphabet{\mathcalligra}{T1}{calligra}{m}{n}
\DeclareFontShape{T1}{calligra}{m}{n}{<->s*[2.2]callig15}{}
\newcommand{\scripty}[1]{\ensuremath{\mathcalligra{#1}}}

% Unit commands
\newcommand{\degree}{\text{°}}
\newcommand{\unitSpeed}{\frac{\text{m}}{\text{s}}}
\newcommand{\unitAcceleration}{\frac{\text{m}}{\text{s}^2}}

% Constant commands
\newcommand{\avogadro}{6.022\times10^{23}}

\newcommand{\boltzmannConstK}{1.381\times10^{-23}\frac{\text{J}}{\text{K}}}
\newcommand{\boltzmannConstInEV}{8.617\times10^{-5}\frac{\text{eV}}{\text{K}}}
\newcommand{\planckConst}{6.626\times10^{-34}\text{J}\cdot\text{s}}

\newcommand{\atmP}{1.01325\times10^5\,\text{Pa}}

% Known Value Commands
\newcommand{\electronMass}{9.109\times10^{-31}\,\text{kg}}
\newcommand{\protonMass}{1.673\times10^{-27}\,\text{kg}}

% Commands to easily write partial derivatives
\newcommand{\partDeriv}[2]{\frac{\partial #1}{\partial #2}}
\newcommand{\doublePartDeriv}[2]{\frac{\partial^2 #1}{\partial^2 #2}}


\title{Problem 7.31}
\author{Gabriel Decker}


\begin{document}


\maketitle
\begin{flushleft}



In this problem, we'll continue our work in problem 7.28 where we found the density of states for a two-dimensional Fermi gas.
We will find the system's total energy $U$ using the Sommerfield Expansion so that we can get the heat capacity $C_V$.
We will begin with the following $U$ equation.
\begin{equation}
    \label{eq:start}
    U=\int_0^\infty\epsilon g(\epsilon)\frac{1}{e^{(\epsilon-\mu)/kT}+1}\,d\epsilon
\end{equation}
with
\begin{equation}
    g(\epsilon)=\frac{4\pi mA}{h^2}=\frac{N}{\epsilon_F}
\end{equation}
to eventually get the heat capacity.
\begin{equation}
    C_V=\left(\partDeriv{U}{T}\right)_V
\end{equation}
Also, recall the following.
\begin{equation}
    \overline{n}_\text{FD}(\epsilon)=\frac{1}{e^{(\epsilon-\mu)/kT}+1}
\end{equation}
and
\begin{equation}
    -\frac{d\overline{n}_\text{FD}(\epsilon)}{d\epsilon}=\frac{1}{kT}\frac{e^x}{(e^x+1)^2}
\end{equation}
with $x=(\epsilon-\mu)/kT$.\\~\\

Plugging in $g$ into $U$, we get the following.
$$U=\int_0^\infty\epsilon g(\epsilon)\frac{1}{e^{(\epsilon-\mu)/kT}+1}\,d\epsilon$$
$$\implies U=\frac{N}{\epsilon_F}\int_0^\infty\epsilon\frac{1}{e^{(\epsilon-\mu)/kT}+1}\,d\epsilon$$
We'll start the Sommerfield expansion by integrating by parts.\\
$dv=\frac{N}{\epsilon_F}\epsilon\,d\epsilon$, so $v=\frac{N}{2\epsilon_F}\epsilon^2$.\\
$u=\overline{n}_\text{FD}(\epsilon)$, so $du=\frac{d\overline{n}_\text{FD}(\epsilon)}{d\epsilon}$.\\
$$\implies U=\frac{N}{2\epsilon_F}\epsilon\overline{n}_\text{FD}(\epsilon)\Big|_0^\infty-\frac{N}{2\epsilon_F}\int_0^\infty\epsilon^2\frac{d\overline{n}_\text{FD}(\epsilon)}{d\epsilon}\,d\epsilon$$
Due to the same logic as problem 7.29 and the book's example, the first term vanishes.
$$\implies U=\frac{N}{2\epsilon_F}\int_0^\infty\epsilon^2\left(-\frac{d\overline{n}_\text{FD}(\epsilon)}{d\epsilon}\right)\,d\epsilon$$
$$\implies U=\frac{N}{2\epsilon_FkT}\int_0^\infty\epsilon^2\frac{e^x}{(e^x+1)^2}\,d\epsilon$$
with $x=(\epsilon-\mu)/kT$ and $dx=d\epsilon/kT$.\\~\\

Changing the integration variable to $x$, $\epsilon=xkT+\mu$ implies the following.
$$\implies U=\frac{N}{2\epsilon_F}\int_{-\mu/kT}^\infty\epsilon^2\frac{e^x}{(e^x+1)^2}\,dx$$
$$\implies U=\frac{N}{2\epsilon_F}\int_{-\mu/kT}^\infty(xkT+\mu)^2\frac{e^x}{(e^x+1)^2}\,dx$$
$$\implies U=\frac{N}{2\epsilon_F}\int_{-\mu/kT}^\infty((kT)^2x^2+2\mu kTx+\mu^2)\frac{e^x}{(e^x+1)^2}\,dx$$
We can also extend the bounds of integration to negative infinity because the extension will add negligibly to the integral.
$$\implies U=\frac{N}{2\epsilon_F}\int_{-\infty}^\infty((kT)^2x^2+2\mu kTx+\mu^2)\frac{e^x}{(e^x+1)^2}\,dx$$
$$\implies U=\frac{N}{2\epsilon_F}\left[\mu^2\int_{-\infty}^\infty\frac{e^x}{(e^x+1)^2}\,dx+2\mu kT\int_{-\infty}^\infty\frac{xe^x}{(e^x+1)^2}\,dx+(kT)^2\int_{-\infty}^\infty\frac{x^2e^x}{(e^x+1)^2}\,dx\right]$$\\~\\

As discussed in the problem 2.29, we have the following results from the book.
$$\int_{-\infty}^\infty\frac{e^x}{(e^x+1)^2}\,dx=1$$
$$\int_{-\infty}^\infty\frac{xe^x}{(e^x+1)^2}\,dx=0$$
$$\int_{-\infty}^\infty\frac{x^2e^x}{(e^x+1)^2}\,dx=\frac{\pi^2}{3}$$\\~\\

Plugging these into our equation for $U$, we get the following.
$$\implies U=\frac{N}{2\epsilon_F}\left[\mu^2\cdot1+2\mu kT\cdot0+(kT)^2\cdot\frac{\pi^2}{3}\right]$$
$$\implies U=\frac{N}{2\epsilon_F}\left[\mu^2+(kT)^2\frac{\pi^2}{3}\right]$$
$$\implies U=\frac{N}{2\epsilon_F}\mu^2+\frac{N}{2\epsilon_F}(kT)^2\frac{\pi^2}{3}$$
Assuming that $\mu\approx\epsilon_F$, we get the following.
$$\implies U=\frac{\epsilon_F}{2}N+\frac{\pi^2}{6}\frac{(kT)^2}{\epsilon_F}N$$
This is a great approximation of the total energy when $kT\ll\epsilon_F$.\\~\\

By applying the equation for $C_V$, we get the following heat capacity for when $kT\ll\epsilon_F$.
$$C_V=\left(\partDeriv{U}{T}\right)_V$$
$$\implies C_V=\partDeriv{}{T}\left(\frac{\epsilon_F}{2}N+\frac{\pi^2}{6}\frac{(kT)^2}{\epsilon_F}N\right)$$
$$\implies C_V=\frac{\pi^2}{3}\frac{kT}{\epsilon_F}Nk$$
$$\implies C_V=\frac{\pi^2Nk^2T}{3\epsilon_F}$$\\~\\

When $kT\gg\epsilon_F$, things are different.
We can go back to the integral of $U$ and $g(\epsilon)$.
$$U=\frac{N}{\epsilon_F}\int_0^\infty\epsilon\frac{1}{e^{(\epsilon-\mu)/kT}+1}\,d\epsilon$$
Recall the result of $\mu$ from part e of problem 7.28, where we found $\mu$ when $kT\gg\epsilon_F$.
$$\mu=kT\ln\left(\frac{\epsilon_F}{kT}\right)$$
Notice that it is negative when $kT\gg\epsilon_F$.
Plugging this definition of $\mu$ into the equation for $U$ yeilds the following.
$$\implies U=\frac{N}{\epsilon_F}\int_0^\infty\epsilon\frac{1}{e^{\epsilon/kT-\ln(\epsilon_F/kT)}+1}\,d\epsilon$$
Since $kT\gg\epsilon_F$, that exponential will always result in a fairly large number, even at $\epsilon=0$.
So, we can treat the $+1$ as negligible.
$$\implies U=\frac{N}{\epsilon_F}\int_0^\infty\epsilon e^{-(\epsilon-\mu)/kT}\,d\epsilon$$\\~\\

This is an integral we can solve (I'm using sympy).\\
$$\implies U=\frac{N}{\epsilon_F}\cdot\left[T^2k^2e^{\mu/kT}\right]$$
Plugging in our $\mu$, we get the following.
$$\implies U=\frac{N}{\epsilon_F}\cdot\left[T^2k^2e^{\ln(\epsilon_F/kT)}\right]$$
$$\implies U=\frac{N}{\epsilon_F}\cdot\left[T^2k^2\frac{\epsilon_F}{kT}\right]$$
$$\implies U=NkT$$
$$\implies C_V=Nk$$
This is the same as an ideal gas with degrees of freedom $f=2$, which makes sense for our two-dimensional problem.




\end{flushleft}

\end{document}
